La première étape de réalisation de ce projet consiste au choix des différents composants le constituant.

Il est important de noter que les critères de choix ne sont pas toujours issus directement du cahier des charges. Cependant, ils permettent une sélection objective des meilleurs composants.

\section{Moteurs}

\subsection{Critères de choix}

\begin{itemize}
    \item \textbf{Vitesse de rotation élevée}  
    Les moteurs doivent pouvoir tourner rapidement afin de permettre une régulation efficace de la position du robot et de le maintenir en équilibre.

    \item \textbf{Précision du contrôle de vitesse}  
    La vitesse des moteurs doit être ajustable avec précision afin de réguler efficacement le balancement du système.  
    Les moteurs doivent également rester précis même à faible vitesse.

    \item \textbf{Synchronisation des moteurs}  
    Les moteurs doivent fonctionner de manière synchrone. Une différence importante de vitesse entre les deux moteurs peut entraîner une rotation non désirée du robot.

    \item \textbf{Consommation énergétique}  
    Les moteurs ne doivent pas consommer trop d’énergie. Étant donné que le système embarqué fonctionne sur batterie, il est important de limiter la consommation afin de préserver l’autonomie.

    \item \textbf{Système de contrôle}  
    Le système de commande des moteurs doit rester simple. Il convient d’éviter une électronique de contrôle trop complexe.
\end{itemize}

\begin{table}[h]
    \centering
    \caption{Synthèse des critères de choix}
    \label{tab:criteres_moteurs}
    \begin{tabular}{|l|p{10cm}|}
        \hline
        \textbf{Critère} & \textbf{Justification} \\
        \hline
        Vitesse de rotation élevée & Permet une réaction rapide pour la stabilisation du robot. \\
        \hline
        Précision du contrôle de vitesse & Assure un contrôle fin du balancement et de la stabilité. \\
        \hline
        Synchronisation des moteurs & Évite les rotations parasites du robot. \\
        \hline
        Consommation énergétique & Améliore l’autonomie du système embarqué sur batterie. \\
        \hline
        Système de contrôle simple & Réduit la complexité électronique et facilite l’implémentation. \\
        \hline
    \end{tabular}
\end{table}